% ══════════════════════════════════════════════════════════════
%  Section 9.9: Transaction to Block (End-to-End Walkthrough)
% ══════════════════════════════════════════════════════════════
\chapter{Chain: Transaction to Block}
\label{ch:tx-to-block}

\begin{learnbox}
\begin{itemize}
  \item Follow the complete path from user input to mined block
  \item Understand how two-node architecture separates web server (API) from miner
  \item Runtime wiring: entry points, message flow, and coordination
  \item The HTTP API entry point and method dispatch
  \item How transactions propagate over P2P and reach miners
  \item Mining triggers, block construction, and broadcast
  \item The complete end-to-end path from \code{main()} to block creation
\end{itemize}
\end{learnbox}

This section is the \textbf{final \index{code walkthrough}code walkthrough} for the \code{\index{chain}chain/} chapter series. The goal is explicit:

\begin{quote}
Follow the first \index{instruction}instruction executed in \code{\index{main}main}, and trace the \index{call chain}call chain all the way to a \index{block}block being created.
\end{quote}

Unlike earlier sections (which focus on individual \index{subsystem}subsystems), this chapter is about \textbf{\index{wiring}wiring and \index{flow}flow}: the \index{runtime}runtime \index{entry point}entrypoints, the \index{message}message boundaries, and the concrete method-to-method path that turns ``a user wants to send \index{bitcoin}bitcoin'' into ``a \index{miner}miner produces a new \index{block}block''.

% ──────────────────────────────────────────────────────────────
\section{A Practical Constraint: Web and Miner Are Separate Roles}
\label{sec:ch09-9-separation}

In \code{\index{main}bitcoin/src/main.rs}, \code{\index{startnode}startnode} enforces that \textbf{a \index{node}node cannot run the \index{web server}web server and be a \index{miner}miner at the same time}. \index{Mining}Mining only runs on \index{node}nodes started with \code{\index{is\_miner}is\_miner=yes}, while \index{HTTP}HTTP \index{transaction}transaction creation requires \code{\index{is\_web\_server}is\_web\_server=yes}.

Therefore, the end-to-end story naturally involves two \index{process}processes:

\begin{itemize}
  \item \textbf{\index{node}Node A (\index{web server}Web/API \index{node}node)}: accepts an \index{HTTP}HTTP \index{request}request and creates a \index{transaction}transaction, then propagates it over \index{peer-to-peer}P2P.
  \item \textbf{\index{node}Node B (\index{miner}Miner \index{node}node)}: receives the \index{transaction}transaction over \index{peer-to-peer}P2P, places it in the \index{mempool}mempool, and \index{mining}mines a \index{block}block once \index{mining}mining is triggered.
\end{itemize}

This is not a \index{limitation}limitation of \index{Bitcoin}Bitcoin as a \index{protocol}protocol---it is simply how this learning \index{implementation}implementation is configured in \code{\index{main}main}.

% ──────────────────────────────────────────────────────────────
\section{Diagram: The Full Pipeline}
\label{sec:ch09-9-pipeline}

\begin{figure}[htbp]
\centering
\begin{tikzpicture}[
  node distance = 0.7cm and 1.0cm,
  web node/.style = {rectangle, rounded corners=5pt, draw=brand, fill=brandlight,
                     text width=2.8cm, align=center, font=\sffamily\small},
  miner node/.style = {rectangle, rounded corners=5pt, draw=sectioncolor, fill=sectioncolor!10,
                       text width=2.8cm, align=center, font=\sffamily\small},
  arrow/.style = {-{Stealth[length=4pt,width=3pt]}, thick, color=brand},
  net arrow/.style = {-{Stealth[length=4pt,width=3pt]}, thick, color=codeframe, dashed},
]

% Node A (Web)
\node[web node] (web_header) at (0, 8.5) {
  \textbf{Node A} \\
  \small Web/API node
};

% Node B (Miner)
\node[miner node] (miner_header) at (6, 8.5) {
  \textbf{Node B} \\
  \small Miner node
};

% Node A steps
\node[web node, minimum height=0.6cm] (http) at (0, 7.5) {
  HTTP POST /tx
};
\node[web node, minimum height=0.6cm] (handler) at (0, 6.5) {
  send\_transaction
};
\node[web node, minimum height=0.6cm] (btc_tx) at (0, 5.5) {
  btc\_transaction
};
\node[web node, minimum height=0.6cm] (process_tx) at (0, 4.5) {
  process\_transaction
};
\node[web node, minimum height=0.6cm] (mempool) at (0, 3.5) {
  add\_to\_memory\_pool
};
\node[web node, minimum height=0.6cm] (send_inv) at (0, 2.5) {
  send\_inv(Tx)
};

% Node B steps
\node[miner node, minimum height=0.6cm] (recv_inv) at (6, 2.5) {
  receive INV
};
\node[miner node, minimum height=0.6cm] (process_stream) at (6, 1.8) {
  process\_stream
};
\node[miner node, minimum height=0.6cm] (recv_tx) at (6, 1.1) {
  receive Tx
};
\node[miner node, minimum height=0.6cm] (mine_trigger) at (6, 0.4) {
  process\_mine\_block
};
\node[miner node, minimum height=0.6cm] (pow) at (6, -0.3) {
  ProofOfWork::run
};

% Node A arrows
\draw[arrow] (http) -- (handler);
\draw[arrow] (handler) -- (btc_tx);
\draw[arrow] (btc_tx) -- (process_tx);
\draw[arrow] (process_tx) -- (mempool);
\draw[arrow] (mempool) -- (send_inv);

% Network arrow
\draw[net arrow] (send_inv) -- (recv_inv)
  node[midway, above, font=\small\sffamily] {P2P};

% Node B arrows
\draw[arrow] (recv_inv) -- (process_stream);
\draw[arrow] (process_stream) -- (recv_tx);
\draw[arrow] (recv_tx) -- (mine_trigger);
\draw[arrow] (mine_trigger) -- (pow);

\end{tikzpicture}
\caption{End-to-End Flow: Transaction from Web API Node to Miner Block Creation}
\label{fig:ch09-9-tx-to-block}
\end{figure}

% ──────────────────────────────────────────────────────────────
\section{Method Index}
\label{sec:ch09-9-index}

We will walk these methods, in the order you reach them at runtime:

\begin{itemize}
  \item \code{bitcoin/src/main.rs}
    \begin{itemize}
      \item \code{main}
      \item \code{process\_command}
      \item \code{start\_node}
    \end{itemize}

  \item \code{bitcoin/src/node/server.rs}
    \begin{itemize}
      \item \code{Server::run\_with\_shutdown}
    \end{itemize}

  \item \code{bitcoin/src/net/net\_processing.rs}
    \begin{itemize}
      \item \code{process\_stream} (only the message arms relevant to tx/block flow)
      \item \code{send\_inv}, \code{send\_get\_data}, \code{send\_tx}, \code{send\_block}
    \end{itemize}

  \item \code{bitcoin/src/web/server.rs} + \code{bitcoin/src/web/routes/api.rs} + \code{bitcoin/src/web/handlers/transaction.rs}
    \begin{itemize}
      \item \code{WebServer::start\_with\_shutdown}
      \item \code{create\_api\_routes} (the POST \code{/transactions} route)
      \item \code{send\_transaction} (HTTP handler)
    \end{itemize}

  \item \code{bitcoin/src/node/context.rs}
    \begin{itemize}
      \item \code{NodeContext::btc\_transaction}
      \item \code{NodeContext::process\_transaction}
      \item \code{NodeContext::submit\_transaction\_for\_mining}
      \item \code{NodeContext::broadcast\_transaction\_to\_nodes}
    \end{itemize}

  \item \code{bitcoin/src/node/txmempool.rs}
    \begin{itemize}
      \item \code{add\_to\_memory\_pool}
    \end{itemize}

  \item \code{bitcoin/src/node/miner.rs}
    \begin{itemize}
      \item \code{should\_trigger\_mining}
      \item \code{prepare\_mining\_utxo}
      \item \code{process\_mine\_block}
    \end{itemize}
\end{itemize}

Throughout, we omit unrelated class/struct definitions and focus only on the methods that form the end-to-end chain.

% ──────────────────────────────────────────────────────────────
\section{Step 1 --- The Process Starts}
\label{sec:ch09-9-step1}

The first executed instructions in the node runtime are in \code{main}. From there, we parse the CLI and route into \code{start\_node(...)}.

\begin{listing}[htbp]
\caption{Process entrypoint (\code{main})}
\label{lst:9-9-1}
\begin{minted}[fontsize=\footnotesize]{rust}
#[tokio::main]
#[deny(unused_must_use)]
async fn main() {
    // 1) Initialize the global tracing subscriber
    // so later components can log.
    initialize_logging();

    // 2) Parse CLI args (subcommands like
    // startnode/createwallet/etc.).
    let opt = Opt::parse();

    // 3) Dispatch the selected command.
    // For end-to-end node runtime, we care about
    // Command::StartNode.
    if let Err(e) = process_command(opt.command).await
    {
        eprintln!("Error: {}", e);
        std::process::exit(1);
    }
}
\end{minted}
\end{listing}

\begin{listing}[htbp]
\caption{CLI dispatch (\code{process\_command})}
\label{lst:9-9-2}
\begin{minted}[fontsize=\footnotesize]{rust}
async fn process_command(command: Command)
    -> Result<()>
{
    match command {
        // ... other commands omitted
        Command::StartNode {
            is_miner,
            is_web_server,
            connect_nodes,
            wlt_mining_addr,
        } => {
            // Validate and normalize mining address
            let validated_addr =
                WalletAddress::validate(
                    wlt_mining_addr)?;

            // Enter the runtime wiring path
            start_node(
                is_miner,
                is_web_server,
                connect_nodes,
                validated_addr
            ).await
        }
        _ => Ok(()),
    }
}
\end{minted}
\end{listing}

\begin{listing}[htbp]
\caption{Wiring the node runtime (\code{start\_node})}
\label{lst:9-9-3}
\begin{minted}[fontsize=\footnotesize]{rust}
async fn start_node(
    is_miner: IsMiner,
    is_web_server: IsWebServer,
    connect_nodes: Vec<ConnectNode>,
    wlt_mining_addr: WalletAddress,
) -> Result<()> {
    // Configure the process-level GLOBAL_CONFIG
    validate_miner_config(
        &wlt_mining_addr,
        &is_miner,
        &is_web_server
    )?;

    // Open or create blockchain
    let blockchain = open_or_create_blockchain(
        &wlt_mining_addr,
        &connect_nodes
    ).await?;

    let node_context = NodeContext::new(blockchain);
    let socket_addr = GLOBAL_CONFIG.get_node_addr();

    // Convert CLI nodes into a set for bootstrapping
    let connect_nodes_set: HashSet<ConnectNode> =
        connect_nodes.into_iter().collect();

    // Create shutdown signal channel
    let (shutdown_tx, _) =
        tokio::sync::broadcast::channel::<()>(1);

    // Start the P2P network server
    let network_server = Server::new(
        node_context.clone());
    let net_shutdown_rx = shutdown_tx.subscribe();
    let network_handle = tokio::spawn(async move {
        network_server
            .run_with_shutdown(
                &socket_addr,
                connect_nodes_set,
                net_shutdown_rx
            )
            .await;
    });

    // ... (continue with optional web server)
}
\end{minted}
\end{listing}

Once the network server is initialized, the function checks whether to launch an optional HTTP web server (for wallet nodes) or handle miner-only mode. Both paths wait for a shutdown signal:

\begin{minted}[fontsize=\footnotesize,linenos=false]{rust}
    match (is_web_server, is_miner) {
        (IsWebServer::Yes, IsMiner::No) => {
            let web_server =
                create_web_server(node_context);
            let web_handle = tokio::spawn(async move {
                let _ = web_server
                    .start_with_shutdown().await;
            });
            // ... (wait for Ctrl+C or task completion)
            Ok(())
        }
        (IsWebServer::Yes, _) => {
            Err(BtcError::InvalidConfiguration(
                "Web server and miner cannot \
                 be enabled at the same time"
                    .to_string(),
            ))
        }
        _ => {
            // ... (wait for Ctrl+C or task completion)
            Ok(())
        }
    }
}
\end{minted}

% ──────────────────────────────────────────────────────────────
\section{Step 2 --- The P2P Runtime Loop}
\label{sec:ch09-9-step2}

Once \code{start\_node} runs, the P2P server is spawned and begins accepting TCP connections. Each accepted stream is handed to the P2P message dispatcher \code{net\_processing::process\_stream}.

\begin{listing}[htbp]
\caption{P2P accept loop (\code{Server::run\_with\_shutdown})}
\label{lst:9-9-4}
\begin{minted}[fontsize=\footnotesize]{rust}
pub async fn run_with_shutdown(
    &self,
    addrs: &SocketAddr,
    connect_nodes: HashSet<ConnectNode>,
    mut shutdown: tokio::sync::broadcast
        ::Receiver<()>,
) {
    // Bind the TCP listener
    let listener = TcpListener::bind(addrs)
        .await
        .expect("TcpListener bind error");

    // Bootstrap: if not central node,
    // send version handshake
    if !addrs.eq(&CENTERAL_NODE) {
        let best_height = self.node_context
            .get_blockchain_height().await
            .expect("Blockchain read error");
        send_version(&CENTERAL_NODE, best_height)
            .await;
    }
    // ... (collect and share known nodes)

    loop {
        tokio::select! {
            _ = shutdown.recv() => break,
            accept_res = listener.accept() => {
                if let Ok((stream, _)) = accept_res {
                    let blockchain = self
                        .node_context.clone();
                    tokio::spawn(async move {
                        if let Ok(std_stream) =
                            stream.into_std()
                        {
                            let _ =
                                net_processing
                                ::process_stream(
                                    blockchain,
                                    std_stream
                                ).await;
                        }
                    });
                }
            }
        }
    }
}
\end{minted}
\end{listing}

% ──────────────────────────────────────────────────────────────
\section{Step 3 --- The HTTP Entrypoint (Node A)}
\label{sec:ch09-9-step3}

The web server is only enabled for non-miner nodes. Its purpose is to provide a client-friendly HTTP API for creating wallets, querying balances, and creating transactions.

\begin{listing}[htbp]
\caption{Web server startup (\code{WebServer::start\_with\_shutdown})}
\label{lst:9-9-5}
\begin{minted}[fontsize=\footnotesize]{rust}
pub async fn start_with_shutdown(
    &self,
) -> Result<(), Box<dyn std::error::Error
                            + Send + Sync>>
{
    let app = self.create_app()?;
    let addr = SocketAddr::from((
        [0, 0, 0, 0],
        self.config.port
    ));
    let listener = tokio::net::TcpListener
        ::bind(addr).await?;
    let shutdown_signal = async {
        tokio::signal::ctrl_c().await?
    };
    axum::serve(
        listener,
        app.into_make_service_with_connect_info::<
            SocketAddr
        >()
    )
        .with_graceful_shutdown(shutdown_signal)
        .await?;
    Ok(())
}
\end{minted}
\end{listing}

\begin{listing}[htbp]
\caption{Route mapping: POST \code{/api/v1/transactions}}
\label{lst:9-9-6}
\begin{minted}[fontsize=\footnotesize]{rust}
pub fn create_api_routes()
    -> Router<Arc<NodeContext>>
{
    Router::new()
        // ... other routes omitted ...
        .route(
            "/transactions",
            post(transaction::send_transaction)
        )
        // ... other routes omitted ...
}
\end{minted}
\end{listing}

\begin{listing}[htbp]
\caption{HTTP handler: \code{send\_transaction}}
\label{lst:9-9-7}
\begin{minted}[fontsize=\footnotesize]{rust}
pub async fn send_transaction(
    State(node): State<Arc<NodeContext>>,
    Json(request): Json<SendTransactionRequest>,
) -> Result<
    Json<ApiResponse<SendBitCoinResponse>>,
    StatusCode
> {
    let txid = node
        .btc_transaction(
            &request.from_address,
            &request.to_address,
            request.amount
        )
        .await
        .map_err(|_| StatusCode::BAD_REQUEST)?;
    let response = SendBitCoinResponse {
        txid,
        timestamp: chrono::Utc::now(),
    };
    Ok(Json(ApiResponse::success(response)))
}
\end{minted}
\end{listing}

% ──────────────────────────────────────────────────────────────
\section{Step 4 --- From Send to Mempool Accept}
\label{sec:ch09-9-step4}

Now we enter the core node orchestration API.

\begin{listing}[htbp]
\caption{Create and submit signed UTXO transaction (\code{NodeContext::btc\_transaction})}
\label{lst:9-9-8}
\begin{minted}[fontsize=\footnotesize]{rust}
pub async fn btc_transaction(
    &self,
    wlt_frm_addr: &WalletAddress,
    wlt_to_addr: &WalletAddress,
    amount: i32,
) -> Result<String> {
    // 1) Build a UTXO view over the current
    // chain state. This is used to select
    // spendable outputs owned by the sender.
    let utxo_set = UTXOSet::new(
        self.blockchain.clone());

    // 2) Construct and sign the transaction.
    // Transaction::new_utxo_transaction selects
    // inputs, creates outputs (including change),
    // and signs the inputs using sender's key.
    let tx = Transaction::new_utxo_transaction(
        wlt_frm_addr,
        wlt_to_addr,
        amount,
        &utxo_set
    ).await?;

    // 3) Submit to mempool + propagation via
    // shared processing entrypoint.
    let addr_from = crate::GLOBAL_CONFIG
        .get_node_addr();
    self.process_transaction(&addr_from, tx).await
}
\end{minted}
\end{listing}

\begin{listing}[htbp]
\caption{Mempool acceptance entrypoint (\code{NodeContext::process\_transaction})}
\label{lst:9-9-9}
\begin{minted}[fontsize=\footnotesize]{rust}
pub async fn process_transaction(
    &self,
    addr_from: &std::net::SocketAddr,
    utxo: Transaction,
) -> Result<String> {
    // 1) Reject duplicates early
    if transaction_exists_in_pool(&utxo) {
        return Err(
            BtcError
            ::TransactionAlreadyExistsInMemoryPool(
                utxo.get_tx_id_hex()
            )
        );
    }

    // 2) Add to the global mempool and mark
    // outputs as ``reserved'' by mempool policy
    add_to_memory_pool(utxo.clone(),
        &self.blockchain).await?;

    // 3) Background work: propagation and
    // (if miner) mining trigger
    let context = self.clone();
    let addr_copy = *addr_from;
    let tx = utxo.clone();
    tokio::spawn(async move {
        let _ = context
            .submit_transaction_for_mining(
                &addr_copy, tx).await;
    });

    // 4) Return txid immediately (accept-to-mempool)
    Ok(utxo.get_tx_id_hex())
}
\end{minted}
\end{listing}

\begin{listing}[htbp]
\caption{Mempool insert (\code{txmempool::add\_to\_memory\_pool})}
\label{lst:9-9-10}
\begin{minted}[fontsize=\footnotesize]{rust}
pub async fn add_to_memory_pool(
    tx: Transaction,
    blockchain_service: &BlockchainService,
) -> Result<()> {
    // 1) Store the transaction in the global
    // in-memory pool.
    GLOBAL_MEMORY_POOL.add(tx.clone())
        .expect("Memory pool add error");

    // 2) Mark referenced outputs as ``in mempool''
    // inside the UTXO view. This is a local
    // double-spend protection mechanism: it prevents
    // two mempool txs from spending the same output
    // concurrently.
    let utxo_set = UTXOSet::new(
        blockchain_service.clone());
    utxo_set.set_global_mem_pool_flag(
        &tx.clone(), true).await?;

    Ok(())
}
\end{minted}
\end{listing}

% ──────────────────────────────────────────────────────────────
\section{Step 5 --- P2P Propagation: INV \texorpdfstring{$\to$}{->} GETDATA \texorpdfstring{$\to$}{->} TX}
\label{sec:ch09-9-step5}

Transactions propagate over the P2P layer, not the HTTP layer. The web layer exists only for local clients (UIs/admin tooling).

\begin{listing}[htbp]
\caption{Background ``after accept'' hook (\code{submit\_transaction\_for\_mining})}
\label{lst:9-9-11}
\begin{minted}[fontsize=\footnotesize]{rust}
async fn submit_transaction_for_mining(
    &self,
    addr_from: &std::net::SocketAddr,
    utxo: Transaction,
) -> Result<()> {
    // Propagate tx to peers if this is central node
    if GLOBAL_CONFIG.get_node_addr()
        .eq(&CENTERAL_NODE)
    {
        let nodes = self
            .get_nodes_excluding_sender(addr_from)
            .await?;
        self.broadcast_transaction_to_nodes(
            &nodes,
            utxo.get_id_bytes()
        ).await;
    }
    // Trigger mining if threshold met
    if should_trigger_mining() {
        if let Some(mining_address) =
            GLOBAL_CONFIG.get_mining_addr()
        {
            if let Ok(txs) =
                prepare_mining_utxo(&mining_address)
            {
                if !txs.is_empty() {
                    process_mine_block(txs,
                        &self.blockchain).await?;
                }
            }
        }
    }
    Ok(())
}
\end{minted}
\end{listing}

\begin{listing}[htbp]
\caption{Broadcast tx inventory to peers (\code{broadcast\_transaction\_to\_nodes})}
\label{lst:9-9-12}
\begin{minted}[fontsize=\footnotesize]{rust}
async fn broadcast_transaction_to_nodes(
    &self,
    nodes: &[Node],
    txid: Vec<u8>
) {
    // For each peer, spawn a task that sends
    // an INV message. INV is a lightweight
    // announcement: ``I have txid''.
    let txid_clone = txid.clone();
    nodes.iter().for_each(|node| {
        let node_addr = node.get_addr();
        let txid = txid_clone.clone();
        tokio::spawn(async move {
            send_inv(&node_addr, OpType::Tx,
                &[txid]).await;
        });
    });
}
\end{minted}
\end{listing}

\begin{listing}[htbp]
\caption{P2P send primitives (\code{send\_inv}, \code{send\_get\_data}, \code{send\_tx}, \code{send\_block})}
\label{lst:9-9-13}
\begin{minted}[fontsize=\footnotesize]{rust}
pub async fn send_get_data(
    addr_to: &SocketAddr,
    op_type: OpType,
    id: &[u8]
) {
    send_data(addr_to, Package::GetData {
        addr_from: GLOBAL_CONFIG.get_node_addr(),
        op_type,
        id: id.to_vec(),
    }).await;
}

pub async fn send_inv(
    addr_to: &SocketAddr,
    op_type: OpType,
    blocks: &[Vec<u8>],
) {
    send_data(addr_to, Package::Inv {
        addr_from: GLOBAL_CONFIG.get_node_addr(),
        op_type,
        items: blocks.to_vec(),
    }).await;
}

pub async fn send_tx(
    addr_to: &SocketAddr,
    tx: &Transaction
) {
    send_data(addr_to, Package::Tx {
        addr_from: GLOBAL_CONFIG.get_node_addr(),
        transaction: tx.serialize()
            .expect("Serialization error"),
    }).await;
}

pub async fn send_block(
    addr_to: &SocketAddr,
    block: &Block
) {
    send_data(addr_to, Package::Block {
        addr_from: GLOBAL_CONFIG.get_node_addr(),
        block: block.serialize()
            .expect("Serialization error"),
    }).await;
}
\end{minted}
\end{listing}

\begin{listing}[htbp]
\caption{P2P receive: miner gets full tx (\code{process\_stream}, tx-related arms)}
\label{lst:9-9-14}
\begin{minted}[fontsize=\footnotesize]{rust}
// Source: bitcoin/src/net/net_processing.rs
match pkg {
    // Peer says: "I have txid"
    Package::Inv {
        addr_from,
        op_type: OpType::Tx,
        items
    } => {
        let txid = items.first()
            .expect("INV items empty");
        let txid_hex = HEXLOWER.encode(txid);

        // If not already in mempool, request
        // the full tx bytes
        if !GLOBAL_MEMORY_POOL
            .contains(txid_hex.as_str())
            .expect("mempool contains error")
        {
            send_get_data(&addr_from, OpType::Tx,
                txid).await;
        }
    }

    // Peer sends: full tx bytes
    Package::Tx { addr_from, transaction } => {
        let tx = Transaction::deserialize(
            transaction.as_slice())
            .expect("Transaction deserialization");

        // Admit to mempool; may also trigger mining
        let _ = node_context
            .process_transaction(&addr_from, tx)
            .await;
    }

    // Peer sends: full block bytes
    Package::Block { addr_from, block } => {
        let block = Block::deserialize(
            block.as_slice())
            .expect("Block deserialization");
        node_context.add_block(&block).await
            .expect("Blockchain write error");
    }

    _ => {}
}
\end{minted}
\end{listing}

% ──────────────────────────────────────────────────────────────
\section{Step 6 --- Mining: Mempool to Block}
\label{sec:ch09-9-step6}

The mining module is the last step in the pipeline. It is only active when the process is started as a miner (\code{GLOBAL\_CONFIG.is\_miner()}).

\begin{listing}[htbp]
\caption{Mining trigger and block assembly (\code{miner::*})}
\label{lst:9-9-15}
\begin{minted}[fontsize=\footnotesize]{rust}
pub fn should_trigger_mining() -> bool {
    let pool_size = GLOBAL_MEMORY_POOL.len()
        .expect("Memory pool length error");
    let is_miner = GLOBAL_CONFIG.is_miner();
    pool_size >= TRANSACTION_THRESHOLD && is_miner
}

pub fn prepare_mining_utxo(
    mining_address: &WalletAddress,
) -> Result<Vec<Transaction>> {
    // Pull all mempool txs and append a
    // coinbase tx for the miner reward
    let txs = GLOBAL_MEMORY_POOL.get_all()?;
    let coinbase_tx =
        Transaction::new_coinbase_tx(mining_address)?;
    let mut final_txs = txs;
    final_txs.push(coinbase_tx);
    Ok(final_txs)
}

pub async fn process_mine_block(
    txs: Vec<Transaction>,
    blockchain: &BlockchainService,
) -> Result<Block> {
    // Construct+mine a new block on the
    // current tip
    let new_block = blockchain.mine_block(&txs)
        .await?;

    // Remove txs from mempool once mined
    for tx in &txs {
        remove_from_memory_pool(
            tx.clone(), blockchain).await;
    }

    // Announce new block hash to peers (INV)
    broadcast_new_block(&new_block).await?;
    Ok(new_block)
}
\end{minted}
\end{listing}

% ──────────────────────────────────────────────────────────────
\subsection{Where the Block Is Actually Created}
\label{sec:ch09-9-where-created}

The exact method call edge where a ``block comes into existence'' is:

\begin{minted}[fontsize=\footnotesize,linenos=false]{text}
miner::process_mine_block
  +-> BlockchainService::mine_block
        +-> BlockchainFileSystem::mine_block
              +-> Block::new_block
                    +-> ProofOfWork::run
\end{minted}

You have already studied these methods in detail in \textbf{Chapter~16 (Consensus and Validation)}; this chapter's purpose is to show \textbf{how we reach them at runtime} from \code{main}.

% ──────────────────────────────────────────────────────────────
\section{Summary}
\label{sec:ch09-9-summary}

\begin{chaptersummary}
  \item The end-to-end flow from user input to block creation spans two separate nodes: a web/API node and a miner node.
  \item \code{main()} parses CLI arguments and dispatches to \code{start\_node}, which initializes the blockchain and spawns P2P and optional web servers.
  \item The HTTP API receives a POST request, constructs a signed UTXO transaction, and submits it to the mempool via \code{NodeContext::process\_transaction}.
  \item Transactions propagate over P2P using INV/GETDATA/TX messages, reaching other nodes including miners.
  \item When the mempool reaches a threshold size and the node is a miner, \code{should\_trigger\_mining} initiates mining.
  \item \code{prepare\_mining\_utxo} snapshots the mempool and appends a coinbase transaction for the block.
  \item \code{Block::new\_block} constructs the header and runs proof-of-work via \code{ProofOfWork::run}.
  \item Once mining completes, the mined block is broadcast via INV/GETDATA/BLOCK messages to peers.
  \item Understanding this end-to-end path clarifies how individual subsystems (transaction construction, signing, mining, consensus) integrate into a complete node runtime.
\end{chaptersummary}
